%! Piecewise & Step Functions — Unit 1, Lesson 1.4 — Homework (Answer Key)
% Mirrors homework/main.tex EXACTLY: every \blank{} becomes an \ans{}, every \writelines{n} is
% answered with exactly n \ansline{}s, and every work block is authored byte-identically, so the
% two files come out the same number of pages. No teacher notes here — they live in the plan.
\documentclass[12pt]{article}
\usepackage{algebra2-article}
\usepackage{algebra2-key}   % pulls in -boxes; do NOT also load -boxes
% Coordinate grid: {xmin}{xmax}{ymin}{ymax}
\newcommand{\lgrid}[4]{%
  \draw[step=1, linegray!40, very thin] (#1,#3) grid (#2,#4);
  \draw[->, semithick, charcoal] (#1,0)--(#2,0) node[below right, font=\scriptsize]{$x$};
  \draw[->, semithick, charcoal] (0,#3)--(0,#4) node[left, font=\scriptsize]{$y$};}
% Endpoint dots: {color}{x}{y}. Closed = filled; open = white-filled ring.
\newcommand{\fdot}[3]{\fill[#1] (#2,#3) circle (3.2pt);}
\newcommand{\hdot}[3]{\draw[very thick, #1, fill=white] (#2,#3) circle (3.2pt);}

\begin{document}

\pageheader{Unit 1, Lesson 1.4}{Homework --- Answer Key}
% NAMESTRIP: no \namedateperiod here — mirrors the blank.

\vspace{0.10in}

\boxguard[8]
\begin{remindbox}
\small
\textbf{This is your graded homework.} Your packet with completed homework is due the first
class after two study halls. I will announce the due date in class and on TurtleNet. Remember:
first find the piece whose condition the input satisfies --- at a \textbf{boundary}, the piece
whose inequality \emph{includes} it --- then apply that rule. A closed dot includes its endpoint;
an open dot excludes it.
\end{remindbox}

\vspace{0.10in}

\boxguard[14]
\begin{notesbox}{Practice}
\small
\begin{enumerate}[leftmargin=1.9em, itemsep=8pt]
  \item \textbf{Evaluate.} $f(x)=\begin{cases} x+4, & x\le -1 \\ 2x, & x>-1 \end{cases}$
    \\[3pt]
    $f(-3)=$ \ans{$1$} \quad $f(-1)=$ \ans{$3$} \quad $f(0)=$ \ans{$0$} \quad
    $f(5)=$ \ans{$10$} \\[3pt]
    Which piece owns $x=-1$, and why? \ans{the first: $-1\le-1$ includes it, and $-1>-1$ is false}

  \item \textbf{Meets or jumps?} Two functions, one number apart:
    \[ g(x)=\begin{cases} x+2, & x<1 \\ 5-x, & x\ge1 \end{cases} \qquad\qquad
       k(x)=\begin{cases} x+2, & x<1 \\ 4-x, & x\ge1 \end{cases} \]
    \vspace{-8pt}
    \begin{center}
    \renewcommand{\arraystretch}{1.4}
    \begin{tabularx}{\linewidth}{@{}X|X@{}}
    \textbf{(a) $g$} & \textbf{(b) $k$} \\
    \hline
    At $x=1$: first rule \ans{$3$}, second rule \ans{$4$} & At $x=1$: first rule \ans{$3$}, second rule \ans{$3$} \\
    So the graph \ans{jumps} (meets / jumps) & So the graph \ans{meets} (meets / jumps) \\
    \end{tabularx}
    \end{center}
    (c) In one sentence: what decides whether two pieces meet at a boundary?
    \ansline{Whether both rules give the \emph{same} value at the boundary --- same value, they meet.}
    \ansline{In (a) the rules give $3$ and $4$, a gap of $1$; in (b) both give $3$, so no gap.}

  \item \textbf{Read the graph, then write its rule.} The graph shows \emph{all} of $f$. \\[2pt]
    \begin{minipage}[c]{0.60\linewidth}
    $f(1)=$ \ans{$1$}; \; $f(2)=$ \ans{$-2$} (read the \ans{closed} dot) \\[3pt]
    Domain \ans{$[-2,5]$}; \; increasing on \ans{$[-2,2)$}; \; constant on \ans{$[2,5]$} \\[3pt]
    Continuous at $x=2$? \ans{no} \; Now write the rule: \\[3pt]
    $f(x)=\begin{cases} \text{\ans{$x$}}, & -2\le x<2 \\[3pt] \text{\ans{$-2$}}, & 2\le x\le5 \end{cases}$
    \end{minipage}\hfill
    \begin{minipage}[c]{0.36\linewidth}\centering
    \begin{tikzpicture}[scale=0.38]
      \lgrid{-2.5}{5.5}{-3}{3}
      \draw[very thick, forest] (-2,-2)--(2,2);
      \fdot{forest}{-2}{-2} \hdot{forest}{2}{2}
      \draw[very thick, forest] (2,-2)--(5,-2);
      \fdot{forest}{2}{-2} \fdot{forest}{5}{-2}
    \end{tikzpicture}
    \end{minipage}
\end{enumerate}
\end{notesbox}

\newpage
\begin{notesbox}{Practice, continued}
\small
\begin{enumerate}[leftmargin=1.9em, itemsep=8pt, start=4]
  \item \textbf{The staircase.} Round \emph{down} each time.
    \\[3pt]
    $\lfloor 4.2\rfloor=$ \ans{$4$} \quad $\lfloor -3.1\rfloor=$ \ans{$-4$} \quad
    $\lfloor 6\rfloor=$ \ans{$6$} \quad $\lfloor -0.5\rfloor=$ \ans{$-1$} \quad
    For every $x$ in $[3,4)$, $\lfloor x\rfloor=$ \ans{$3$} \\[3pt]
    For which inputs is $\lfloor x\rfloor=x$ exactly? \ans{the integers} \; Is $\lfloor x\rfloor$
    \emph{ever} bigger than $x$? \ans{no}

  \item \textbf{Model it.} A phone plan costs \$30 for up to $5$ GB of data, then \$10 for each
        GB over $5$. Let $C(g)$ be the monthly cost for $g$ GB.
    \begin{enumerate}[label=(\alph*), leftmargin=1.9em, itemsep=4pt]
      \item Complete the rule:
            $C(g)=\begin{cases} \text{\ans{$30$}}, & 0\le g\le5 \\[3pt] 30+\text{\ans{$10(g-5)$}}, & g>5 \end{cases}$
      \item Find $C(8)$, and say what it means.
        \begin{work}
          C(8) &= 30 + 10(8-5) \\
               &= 60
        \end{work}
        $C(8)$ means: \ans{$8$ GB in a month costs \$60 --- \$30 plus \$10 for each of the $3$ extra GB}
      \item Now the company adds a \$10 \emph{surcharge} the moment you pass $5$ GB. Rewrite the
            second piece: $C(g)=$ \ans{$40+10(g-5)$} for $g>5$. Do the pieces still meet at $g=5$?
            \ans{no} \; The graph now \ans{jumps} from \$\ans{$30$} to \$\ans{$40$}.
    \end{enumerate}

  \item \textbf{Multiple choice (SOL-style).} For
        $f(x)=\begin{cases} 2x+1, & x<2 \\ x-4, & x\ge2 \end{cases}$, what is $f(2)$?
    \begin{enumerate}[label=\Alph*., leftmargin=2.2em, itemsep=1pt]
      \item $5$
      \item \ans{$-2$} \hfill \textcolor{keyred}{\textbf{$\leftarrow$ correct}}
      \item $3$
      \item $f(2)$ is undefined
    \end{enumerate}
    Say in one sentence how you ruled out one of the others.
    \ansline{$x=2$ satisfies $x\ge2$, so the second piece owns it: $f(2)=2-4=-2$.}
    \ansline{A is $2(2)+1$ from the strict-$<$ piece, which does not include $2$; D is wrong because one piece does.}
\end{enumerate}
\end{notesbox}

\vspace{0.10in}

\boxguard[8]
\begin{spiralbox}
\small
\textbf{Coming next --- Lesson 1.5: Linear Regression.} We leave exact rules behind. Given a
\emph{scatter} of real data points that lie on no single line, you will read one \textbf{line of
best fit}, measure how well it fits with the \textbf{correlation coefficient}, and use it to
predict --- the first time in Unit~1 the data is messy.
\end{spiralbox}

\end{document}
